Algorithm~\ref{alg:tt_generative} describes a family of generative optimization
procedures based on tensor-train (TT) probability models. The method maintains
a distribution $P_\theta(x)$ for an integer vector $x=(x_1,\ldots,x_N)$:
\begin{equation}
\begin{gathered}
\pi_\theta(x)=\sum_{\alpha_1,\ldots,\alpha_{N-1}}
G^{(1)}_{1,x_1,\alpha_1}G^{(2)}_{\alpha_1,x_2,\alpha_2}\cdots
G^{(N)}_{\alpha_{N-1},x_N,1},\\
P_\theta(x)=|\pi_\theta(x)|^2/Z_\theta
\quad\text{or}\quad P_\theta(x)=\pi_\theta(x).
\label{eq:tt-def}
\end{gathered}
\end{equation}
The first choice is a Born machine; the second requires a nonnegative,
normalized TT. The implementations considered below use the Born rule.
An initial feasible dataset supplies a starting model and a best-so-far
solution. At each iteration, the model generates candidates, a selection
policy constructs the next training set, and a weighted likelihood update
biases the model toward promising solutions. The incumbent is maintained
separately from training data and is never discarded by diversification.

\textbf{Selection and memory.} The candidate pool may contain only current
TT samples, or also externally generated feasible points and retained data.
The selection policy may favor low cost, charge diversity, or their combination;
it also specifies whether initial points or a history participate. Neither
elite-only selection nor cumulative retention is intrinsic to the framework.
For finite computational cost, experiments specify a selection quota and a
memory policy explicitly.

\textbf{Optional diversification.} Before fitting the selected data, an
operator $\mathcal D_t$ may leave the TT unchanged, reset its amplitudes,
rebuild its charge graph from selected feasible points, or apply a preliminary
sweep with equal sample weights. The last choice flattens the empirical
training target; a single sweep does not guarantee a uniform model over its
entire feasible support. Diversification can be periodic or triggered by
stagnation. It must preserve feasibility and provide nonzero support for
positive-weight training points. These options define different algorithms,
not interchangeable implementations of the same experiment.
With $\mathcal D_t$ equal to the identity, learned amplitudes persist across
iterations. The EELS comparison in \cref{subsec:eels_optimization} uses this
choice; ADD rebuilds the model after selection, while the SA comparison in
\cref{subsec:sa_comparison} uses the explicitly stated reset/rebuild variant.
The schematic loop can be cyclically organized as train--sample--select
when the initial dataset is trained before the first sampling phase.
Detailed training conventions are given in \cref{app:sec:implement_vars}.

